\documentclass{beamer}
\usepackage[utf8]{inputenc}
\usepackage[T1]{fontenc}
\usetheme{Madrid}
\usecolortheme{default}

\title{Designing Markets}
\subtitle{Chapter 4: Handbook of Market Design}
\author{Summary Presentation}
\date{\today}

\begin{document}

\frame{\titlepage}

\begin{frame}{Overview: Three Goals of Market Design}
\begin{enumerate}
    \item \textbf{Diagnose} potential market failures and identify their causes
    \item \textbf{Evaluate} markets and compare alternative designs (qualitative and quantitative effects)
    \item \textbf{Propose} new, improved designs building on lessons learned
\end{enumerate}

\vspace{0.5cm}
\textit{Key insight:} Seemingly small details about market rules can have large impacts on market performance
\end{frame}

\begin{frame}{4.1 Diagnosing Market Failures}
\textbf{Common types of failures documented:}
\begin{itemize}
    \item \textbf{Market Power:} Exercise by large participants (US Spectrum: 7-20\% higher payouts)
    \item \textbf{Collusion:} Centralized mechanisms viewed with suspicion (medical residency lawsuit)
    \item \textbf{Rent Seeking:} Wasteful competition for rents (HFT, ticket brokers)
    \item \textbf{Participation/Entry:} Adding one bidder > optimal design (Bulow-Klemperer)
    \item \textbf{Fragmentation:} Kidney exchange fragmentation reduces transplants 30\%+
    \item \textbf{Search Costs:} Students applying to too few/many schools
    \item \textbf{Thin markets:} Low participation → hard to match
    \item \textbf{Congestion:} Too many participants/transactions → delays/costs
    \item \textbf{Information frictions:} Hidden quality, hidden preferences
    \item \textbf{Market Power:} Strategic Behavior, hidden preferences
    \item \textbf{Strategic Misreporting:} Manipulable mechanisms $\rightarrow$ inefficiency
\end{itemize}
\end{frame}

\begin{frame}{4.2 Evaluating and Comparing Designs}
\textbf{Core Challenge:} Theory often ambiguous—need empirical quantification of trade-offs

\vspace{0.3cm}
\textbf{School Choice (NYC Reform):}
\begin{itemize}
    \item Centralization achieved 80\% of potential gains (old: $\leq 33\%$)
    \item \textbf{Key insight:} Coordination matters more than specific mechanism choice
    \item DA vs. IA: Similar welfare if equilibrium play, but mistakes favor DA
\end{itemize}

\vspace{0.3cm}
\textbf{Dynamic Assignment (organs, housing, licenses):}
\begin{itemize}
    \item Challenge: Objects arrive over time, can't elicit full preferences
    \item FCFS vs. LCFS depends on preference structure (idiosyncratic vs. vertical)
    \item Approach: Dynamic discrete choice models + administrative data
\end{itemize}
\end{frame}

\begin{frame}{4.2 Methodological Approaches}
\textbf{Beyond Structural Models:}
\begin{itemize}
    \item \textbf{Direct comparison:} Two designs in same market (Niederle-Roth: gastro fellowship)
    \item \textbf{Survey-based:} Elicit truthful preferences, simulate (Budish-Cantillon)
    \item \textbf{Descriptive:} Document post-reform changes (Prendergast: food banks)
\end{itemize}

\vspace{0.3cm}
\textbf{Non-Utilitarian Objectives:}
\begin{itemize}
    \item Public housing: targeting vs. efficiency; Organs: survival vs. urgency
    \item Distributional concerns and meritocratic admissions
    \item \textbf{Challenge:} Selection bias when choices correlate with outcomes
    \item \textbf{Solution:} Combine choice models with treatment effect estimation
\end{itemize}

\vspace{0.3cm}
\textit{Note:} Centralized markets provide ideal data for policy evaluation
\end{frame}

\begin{frame}{4.3 Proposing New Designs: Medical Match \& School Choice}
\textbf{Medical Match (NRMP):}
\begin{itemize}
    \item Matches ~25,000 grads annually; discovered Gale-Shapley in 1950s
    \item Solved market unraveling; Roth-Peranson: handled married couples
\end{itemize}

\vspace{0.3cm}
\textbf{School Choice (Abdulkadiroğlu-Sönmez 2003):}
\begin{itemize}
    \item Documented Boston mechanism's manipulability
    \item Proposed DA and TTC as alternatives
    \item DA widely adopted; efficiency differences small but complexity/fairness matter
\end{itemize}
\end{frame}

\begin{frame}{4.3 Proposing New Designs: Course Allocation \& Other Markets}
\textbf{A-CEEI (Budish):}
\begin{itemize}
    \item Allocate oversubscribed courses efficiently and fairly
    \item Approximately efficient, strategy-proof; satisfies envy-freeness
    \item Multi-paper program: theory, computation, experiments, implementation
\end{itemize}

\vspace{0.3cm}
\textbf{Food Banks (Prendergast):}
\begin{itemize}
    \item Pseudomarket with internal currency → equitable, efficient sorting
\end{itemize}

\vspace{0.3cm}
\textbf{Financial Exchanges (Budish et al.):}
\begin{itemize}
    \item Frequent batch auctions eliminate speed arbitrage
    \item Could reduce costs 17\% (~\$5B/year); trade on price not speed
\end{itemize}
\end{frame}

\end{document}
